\section*{第 \ref{ch:b3:perturbation-theory} 章——微扰理论}

\begin{solution}{exo:b3:perturbation-theory:1}
(a) 每个 $|\varphi_n|^2$ 关于 $a/2$ 对称，故 $\langle
x\rangle = a/2$，且 $\Delta E_n = \lambda/2$ 对所有 $n$ 成立。(b)
对称性，而非动力学。(c) 公共平移在每一个差值中都相互抵消。
光谱学测量的是间距，在一级近似下看不到任何变化。
(d) 当 $\lambda$ 与能级间距相当时，
被略去的态混合（二级修正）便不可忽略。
\end{solution}

\begin{solution}{exo:b3:perturbation-theory:2}
$\bra{n\pm1}\hat x\ket n = \sqrt{\hbar/2m\omega}\,\sqrt{n + \tfrac12
\pm \tfrac12}$；两个二级项之和为
\[
\frac{q^2\mathcal E^2\hbar}{2m\omega}\Big[\frac{n+1}{-\hbar\omega}
+ \frac{n}{\hbar\omega}\Big] = -\frac{q^2\mathcal E^2}{2m\omega^2} ,
\]
与 $n$ 无关，恰为精确结果。态\emph{确实}
在一级被修正（$n \pm 1$ 态的混入使波函数
横向移位）；只有能量的一级项
由宇称而消失。
\end{solution}

\begin{solution}{exo:b3:perturbation-theory:3}
(a) $E_\pm = \bar E \pm \sqrt{(\Delta/2)^2 + v^2}$，其中 $\Delta =
E_2 - E_1$。(b) $E_- \approx E_1 - v^2/\Delta$，$E_+ \approx E_2 +
v^2/\Delta$：二级公式逐项成立。(c) 能隙
$2\sqrt{(\Delta/2)^2 + v^2} \ge \Delta$，仅在 $v =
0$ 时取等。(d) 展开式发散，精确解给出 $\pm|v|$——
这正是对简并块中的 $\hat V$ 作对角化
（\cref{thm:b3:perturbation-theory:degenerate}）所规定的做法。
\end{solution}

\begin{solution}{exo:b3:perturbation-theory:4}
(a) 展开式给 $\ket n$ 的系数取 $1$；归一化
一级时迫使修正项落入正交
补空间。(b) 由 $\bra m$ 上的投影得到；小混合要求
$|V_{mn}| \ll |E_n - E_m|$。(c) $(0.1)^2 = 1\%$。(d)
$(0.1/0.01)^2 = 100$：荒谬，应对二能级块
作精确对角化。
\end{solution}

\begin{solution}{exo:b3:perturbation-theory:5}
(a) $\hat x^3$ 为奇：其在任一宇称本征态中的平均值
为零。(b) 将 $\hat x^4 \propto (\hat a + \hat
a^\dagger)^4$，保留升、降次数相等的六项，
得 $\langle(\hat a + \hat a^\dagger)^4\rangle_n = 6n^2
+ 6n + 3$：即所述位移。(c) 当 $\gamma < 0$ 时，位移
按 $n^2$ 的规律变得更负：相邻间距 $E_{n+1} - E_n$
随 $n$ 线性减小。(d) 真实振动能级的收敛梯级
——泛音略小于基音的整数倍，
并有有限的解离极限。
\end{solution}

\begin{solution}{exo:b3:perturbation-theory:6}
(a) 该对态简并，分母为零。(b) 对角元
为 $\lambda(a/2)^2$；非对角元为 $\lambda c^2$，其中 $c =
-8a^2/9\pi^2$——积分分解为两个
一维积分的乘积。(c) 能级位移为 $\lambda[(a/2)^2 \pm
(8a^2/9\pi^2)^2/a^2\cdot a^2]$——即 $\lambda a^2(\tfrac14 \pm
64/81\pi^4)$——本征态为 $(\ket{1,2} \pm \ket{2,1})/
\sqrt2$。(d) 对称性早已给出了本征态（
交换偶与交换奇的组合）；微扰只能
从中选出这些，它确实这样做了。
\end{solution}

\begin{solution}{exo:b3:perturbation-theory:7}
(a) $3ea_0\mathcal E = 3 \times \num{5.29e-11} \times \num{2.5e6}
= \qty{0.40}{meV}$，两侧各为此值。(b) $\bra{100}\hat z\ket{100} = 0$，
由宇称可知，且不存在可与之杂化的简并伙伴。(c)
$\alpha \sim 2e^2a_0^2/E_{\text{I}} \to \alpha/4\pi\varepsilon_0
\sim 2a_0^3$，与精确值 $4.5\,a_0^3$ 同量级。(d)
介电常数：每个电容器的 $\varepsilon_{\text{r}}$
即原子通过恰好这种二级极化来响应电场。
\end{solution}

\begin{solution}{exo:b3:perturbation-theory:8}
(a) $A = \omega^3|d|^2/3\pi\varepsilon_0\hbar c^3 \approx
\qty{6e8}{s^{-1}}$，$\tau = \qty{1.6}{ns}$。(b) $\omega$ 小了
$4.8$ 倍，$\omega^3$ 小 $110$ 倍；偶极矩略大，D
线算得约 $\sim\qty{16}{ns}$——与测量相符。(c)
$1/A = \qty{3.4e14}{s} \approx 11$ 百万年。(d) $\omega^3$
跨越 $10^{20}$ 倍，自紫外至无线电：快速紫外闪光，
地质尺度的无线电耐心——故 21 厘米线
由理论预言（范德胡尔斯特，1944），并在天空中搜寻，
其中 $10^{66}$ 个原子以数量抵消了这份漫长等待。
\end{solution}

\begin{solution}{exo:b3:perturbation-theory:9}
(a) 两态皆偶，与奇的 $\hat z$ 相乘后被积函数为奇。(b)
$\varphi$ 积分 $\int_0^{2\pi}\eu^{\iu m\varphi}\dd\varphi$
除非 $m = 0$，否则为零。(c) 光子携带一个 $\hbar$，且宇称为
奇：$l$ 须改变一，$m$ 至多改变一。(d) 在
所有单光子通道均被关闭时，$2s$ 态寿命为 $\sim\qty{0.12}{s}$——
比 $2p$ 长八个数量级，通过同时发射
\emph{两}个光子衰变，一种微弱的连续谱，在
行星状星云中确已观测到。
\end{solution}

\begin{solution}{exo:b3:perturbation-theory:10}
(a) 将分母替换为最小者 $\tfrac34
E_{\text{I}}$，并利用完备关系：$\alpha/4\pi\varepsilon_0 \le
2\,a_0^2\,(2E_{\text{I}})\,/\,(\tfrac34 E_{\text{I}})\cdot\tfrac12
= \tfrac{16}3\,a_0^3$。(b) 精确值 $4.5\,a_0^3$ 位于
上界之下，理应如此。(c) $n - 1 = N\alpha_{\text{SI}}/2\varepsilon_0
\approx \num{e-4}$，对比实测的 \num{1.3e-4}：原子的
二级软性\emph{正是}气体的折射。(d)
一个矩阵元求和决定了原子如何响应电场，而
每立方米 $10^{25}$ 个原子的集体极化使
光线弯折：\omterm{thm:b3:perturbation-theory:corrections}{微扰理论} $\to$ 极化率 $\to$ 光学。
\end{solution}

\begin{solution}{exo:b3:perturbation-theory:11}
(a) 两支双曲线，最近接近距离 $2|v|$ 位于 $t = 0$ 时。
(b) 在沿下方曲线连续演化的态中——初始的
“非绝热”态的性质已互换：完全
转移。(c) 穿越时间 $\sim v/\lambda$ 远长于
$\hbar/v$，故当 $v^2/\hbar\lambda \gg 1$ 时为绝热过程。(d) 量子比特：
缓慢扫描控制场穿过避免交叉，则布居
沿下方分支从一个基态转移到
另一个基态——这一转移对计时误差稳健，在绝热态制备中
被广泛采用。
\end{solution}

\begin{solution}{exo:b3:perturbation-theory:12}
(a) 求和的每一项都是某个平方除以负的分母。
(b) 两个基态原子之间的主导相互作用
是其偶极–偶极耦合的二阶效应：为负——总是
吸引，对任意一对原子种类皆然。(c) $\bra{\psi_0}\hat H\ket{
\psi_0} = E_0 + \bra{\psi_0}\hat V\ket{\psi_0}$ 是试探态的能量，
因此 $\ge$ 真实基态能量：一阶近似
只能偏高。(d) 斯塔克分裂后的 $n = 2$ 态是激发态；
定理针对的是真实基态，其线性
位移为零——并无冲突。
\end{solution}

\begin{solution}{pb:b3:perturbation-theory:1}
\textbf{1.} $\hat V = e\mathcal E\hat z$ 的一阶平均值为零
（由宇称）；二阶求和中，基态的所有分母均为负，
给出 $-\tfrac12\alpha\mathcal E^2$，其中 $\alpha$ 为
极化率——为负，是因为能级受来自上方的排斥。
\textbf{2.} $\alpha_{\text{SI}} = 4\pi\varepsilon_0 \times
\num{2.2e-30} = \qty{2.4e-40}{F.m^2}$: $p = \qty{2e-37}{C.m}
\approx \num{2e-8}\,ea_0$——约为原子偶极矩的
一亿分之一，按每个分子计。
\textbf{3.} $\hbar\omega \approx \qty{2.3}{eV}$，相比
紫外能隙 $\gtrsim\qty{10}{eV}$：微扰
分母几乎感受不到驱动，因此静态 $\alpha$
可以适用。
\textbf{4.} $p_0 = \alpha\mathcal E_0$ 与频率无关，故 $P
\propto \omega^4$：瑞利定律。
\textbf{5.} $(700/450)^4 = 5.9$：蓝光的散射
强度约为红光的六倍。
\textbf{6.} 太阳发出的紫光少于蓝光，而人眼
对紫光的敏感度很低：散射天空的峰值，\emph{对
我们而言}，位于蓝色。
\textbf{7.} $[\alpha/4\pi\varepsilon_0] = \unit{m^3}$，且
$\omega^4/c^4 = \unit{m^{-4}}$：面积。在 \qty{550}{nm} 处：
$\sigma \approx \qty{7e-31}{m^2}$。
\textbf{8.} $\ell = 1/N\sigma \approx \qty{58}{km}$：相当于
数个大气层厚——空气\emph{几乎}透明。
\textbf{9.} 蓝光（$\ell \approx \qty{26}{km}$）：$1 -
\eu^{-8/26} = 26\%$ 被散射；红光（$\ell \approx \qty{150}{km}$）：
$5\%$。天空的颜色正源于这一差异。
\textbf{10.} 经过约 $\sim\qty{240}{km}$ 的低太阳高度路径后，蓝光
残留 $\eu^{-9} \sim \num{e-4}$，而红光保留 $\sim20\%$：太阳
落时呈红色——同样的红光经地球周缘的日落环折射，
把月食中的月球染成铜色。
\textbf{11.} 偶极子沿其自身轴线不辐射：在
$90^\circ$ 散射时，分子振动
沿视线方向的分量无贡献，而残存的
光偏振方向垂直于太阳–分子–眼睛所在平面。
\textbf{12.} 云层引起多次散射，打乱了
几何关系：偏振指南针在阴天失效——
这恰恰使其所谓的维京人使用成为一项壮举。
\textbf{13.} 关键在于分子的\emph{位置}：空气是随机
气体，其密度在各尺度上涨落——完全有序的物质
（理想情况下为晶体）只会向
折射光束散射。
\textbf{14.} 随机分布、间距为波长量级的散射体，
以随机相位叠加：交叉项平均掉，强度相加
——相消的“共谋”需要有序，无序则将其
废止。
\textbf{15.} $1 - \eu^{-8/58} \approx 13\%$ 的阳光用于
形成蓝色天穹——这与比太阳圆盘暗数千倍、
却仍亮到足以阅读的天空相符。
\textbf{16.} 在液滴内，子波直接叠加\emph{振幅}：
功率 $\propto N^2$，压倒性且几乎与波长无关
（液滴远大于 $\lambda$）：云散射一切光
——呈白色。
\textbf{17.} 亚波长散射体选择颜色
（$\omega^4$）；超波长散射体则按几何光学反射
并呈白色：牛奶、雾、涂料、雪。
\textbf{18.} 接近临界点时，密度涨落
本身增长到光学尺度：流体成为自身的“云”
——临界乳光，即瑞利机制被放大到
不透明的程度。
\textbf{19.} $n - 1 = \num{2.5e25} \times \num{2.4e-40}/(2 \times
\num{8.85e-12}) \approx \num{3.4e-4}$，与实测值
\num{2.9e-4} 相比：前向子波与光束相干，减慢其
相位。
\textbf{20.} 能级的斯塔克位移；天空的蓝色；
光的折射（以及海市蜃楼和空气透镜）——同一个
$\alpha$ 对应三种现象。
\textbf{21.} 臭氧的吸收使天顶暮光呈蓝色；
气溶胶使地平线泛白；多次散射光使
日落后的天空保持微弱发光。
\textbf{22.} 无空气则无偶极子：太阳在黑色天空中耀眼
——以缺席证明，我们的蓝色天穹是散射的阳光，
而非“太空”的一种发光性质。
\textbf{23.} 火星尘埃颗粒相对于 $\lambda$ 较大：
它们在各个方向上高效散射红光（使天空
染色），而蓝光最强烈地\emph{向前}衍射——因此
天空呈奶油糖色，落日周围的光晕为
蓝色：颗粒尺寸颠倒了第 18 题的逻辑。
\textbf{24.} 汤姆孙散射与频率无关：没有颜色
选择性，也无蓝色——但仅与
\emph{电子密度}成正比的响应，正是晶体学与
射线成像所需探针的特性。
\textbf{25.} 一个二阶常数 $\qty{2.2e-30}{m^3}$，
平方并以 $\omega^4$ 加权，得到约 $\sim\qty{60}{km}$
绿光的散射长度：足以在正午保持透明，
足以形成蓝色天穹，其偏振如同一群
受迫偶极子，如期使每次日落泛红。
\end{solution}
